\section{Introduction}

openEMS~\cite{liebig2013} is a widely used free finite-difference time-domain (FDTD)~\cite{yee1966,taflove} solver.
The author runs it from a Python toolkit for printed-circuit-board RF work: 50\,$\Omega$ grounded-coplanar waveguide
(GCPW) lines with lossy conducting sheets, vertical and edge-mount SMA (SubMiniature version A) connector launches with
a coaxial port and a perfectly matched layer (PML), coupled lines and cavities. The machine is an Apple M5 laptop (4
performance and 6 efficiency cores, 16\,GB of unified memory on a single bus of about 150\,GB/s~\cite{m5}, shared with
the GPU). A survey of the toolkit's runs on it found 0.1--2.4 million cells, 41\,000--165\,000 time steps and 10--50
minutes per run at 117--460 million cell updates per second (\MC; from the project log, not in \code{results/}). The
two boards benchmarked here take 76.5\,s and 562.7\,s per full run with stock openEMS (measured,
Table~\ref{tab:e2e}). A parameter study of a launch footprint is a day.

The question this work set out to answer was simply ``how fast can openEMS go on this machine, without changing its
answers?'' The constraint matters: the solver is used as a design gate, its results are compared to measurements and
to earlier runs, and a fast engine that moves the answer by a tenth of a decibel is a liability, not an asset. Every
change to the solver reported here therefore had to reproduce the stock engine's port voltages and currents,
S-parameters and field dumps to the bit for the same number of time steps (the two exceptions, the GPU engines, as
Apple's GPU flushes subnormal numbers to zero, and a float accumulation order in the far-field transform, are called
out where they appear; the fork's own Metal engine is held to the bit against the CPU engine run with the same flush).
These changes make up a fork of openEMS, which together with the benchmark harness and the verification method used
to build it is called fdtdKit (emsKit until 2026-11-05, the name its repository keeps; Section~\ref{sec:repro}).
fdtdKit is the FDTD kit of EMSKit, a family of electromagnetic solvers for Apple silicon; the family's other kits
(femKit, momKit and sbrKit) are not part of this report.

\paragraph{Contributions.}
\begin{itemize}
\item A benchmark harness that runs the stock and the modified engines in-process on the same models with a fixed
  number of time steps, waits for a quiet machine, and compares port probes and S-parameters bit for bit; HDF5 field
  dumps are compared dataset by dataset with \code{scripts/checks/h5cmp.py} (Section~\ref{sec:method}).
\item A sequence of engine changes, each measured and verified on six models (Section~\ref{sec:engine}): thread
  barriers and a per-extension hook plan, a rewritten inner loop, work stealing and an automatic thread count, vector
  paths for the Mur and uniaxial-PML (UPML) boundaries, a fused E/H sweep that lowers the modelled DRAM traffic from
  74 to 49 bytes per cell, the SSE engines' use of upstream's precomputed dump gather, staggered allocations that
  avoid an allocation-placement valley, a scalar-SSE path and compressed tables for PML slabs, line hooks for the
  dispersive-material extension, packed coefficient tables, and thread quality-of-service on macOS. Together
  (snapshot cg14, against stock v0.37.0 built with Apple clang): 3.2--6.9$\times$ on the four models of up to about a
  million cells, 1.7--2.5$\times$ on the two of 5--8 million cells, where stock already ran at up to half the bus
  (Section~\ref{sec:background}). On the patch tutorial about 1.13$\times$ of the 6.9$\times$ comes from the fork's
  base, upstream's \code{gpu-engine} branch, whose HDF5 writer writes that model's time-domain dumps in the
  background; on the notch tutorial, the GCPW line and patch x5 the base runs within 0.98--1.03$\times$ of stock
  (Section~\ref{sec:results}). Snapshot cg18 (patches 0001--0006, the build installed on the author's Mac from
  2026-11-04 to 2026-11-06) gives 3.1--6.8$\times$ on the same four models against the stock rows of its check
  tables. Snapshot cg29 (patches 0001--0027, the installed build from 2026-11-06 to 2026-11-11) adds the CPU items
  of a second
  review and ideas from other contributors' work (Section~\ref{sec:cg29}): an operator index per z-vector, neighbour
  flags in place of the two barriers per time step, a y-row path for UPML boxes, merged ADE updates for conducting
  sheets, geometry coefficient tables and line hooks for the steady-state detection and Debye materials. Paired with
  stock in the same rounds (cg29 at 8 threads, its best setting; stock at its best of 2, 4 and 8 threads), it runs
  the four models 4.73--6.89$\times$ and the two large ones 1.95$\times$ and 2.72$\times$ as fast as stock
  (Table~\ref{tab:current}); paired with cg18 in rounds of its own (5 rounds, other processes on 0.3--0.6 cores; the
  SMA launch on its 0.871\,M-cell mesh), 1.03--1.23$\times$ as fast at 4 threads and 1.13--1.44$\times$ at 8 over the
  six models.
\item Setup and post-processing: a threaded operator construction (4.1\,s to 0.9\,s on the 7.9-million-cell patch
  model), a far-field transform written as complex matrix products, one per plane and block of directions (17.9\,s to
  1.0\,s on the same model; this path needs Apple's Accelerate, other builds keep the original loop), and two small
  CSXCAD changes (times from the project log, not in \code{results/}).
\item An analysis of the ceiling (Section~\ref{sec:beyond}). On a $201^3$ box the CPU engine (snapshot cg9, 8 threads),
  upstream's Metal engine from its pre-merge \code{gpu-engine} branch, and a standalone Swift/Metal kernel that runs
  the exported openEMS operator bit-identically reached 1\,562, 1\,638 and 1\,670\,\MC\ (from the project log, not in
  \code{results/}). By the traffic model the two Metal kernels move 74\,B per cell (separate E and H passes) and the
  fused CPU sweep 49\,B, so the GPU drew more of the bus for about the same cell rate; no byte count was measured with
  hardware counters. On all six test models the CPU engine of cg14 and of cg18 is faster than that Metal engine
  (Section~\ref{sec:ceiling}). Upstream master's merged Metal engine and luikore's diamond-blocked one were built and
  run here on 2026-11-04 (Section~\ref{sec:positioning}), and on 2026-11-11 next to cg31 on luikore's benchmark model
  (Table~\ref{tab:threebuilds}).
\item A diamond-blocked Metal engine (\code{--engine=metal-diamond}, opt-in, Apple only; patches 0007--0016,
  Section~\ref{sec:gpu}) that advances tiles of the grid several time steps at a time, with the tile schedule of
  luikore's openEMS pull request \#228~\cite{luikore228} (none of its code is used). All six models run on the GPU,
  identical bit for bit to the multithreaded CPU engine run with \code{OPENEMS\_FTZ=1}; against stock, which keeps
  subnormal numbers, its port probes differ by $1.7\times10^{-7}$ to $9.4\times10^{-7}$ relative. Paired with the CPU
  engine of cg29 at its best (Table~\ref{tab:current}) it runs patch x5 and notch x3 1.47$\times$ and 1.46$\times$ as
  fast and the GCPW line 1.17$\times$, is level on the SMA launch (1.00$\times$, 1.071\,M cells) and slower on the
  notch and patch tutorials (0.97$\times$ and 0.93$\times$). Across grid sizes (Figure~\ref{fig:crossover}) it leads
  in every round from 1.7--2.9 million cells on, depending on the model, and by 1.25--1.52$\times$ above 3 million; on
  smaller grids it waits for the CPU between short calls (the patch tutorial) or runs at a lower rate of its own (the
  notch tutorial), and the CPU engine is at its fastest between 0.6 and 1.9 million cells. Unblocked, as a fused sweep ($K=1$), the Metal engine runs the two large
  models at 1\,730--1\,840\,\MC, only 67--75\,\% of that sweep's traffic bound (cg27, medians of paired rounds with
  2.5--4 cores of other load; \code{results/fused\_sweep\_k1\_2026-11-04.txt}): the blocking carries the gain.
  A later build, cg31 (patches 0028--0039 of 2026-11-11, the installed build since that day), raises the Metal
  engine's rates on the six models by 1.07--1.53$\times$, which puts it ahead of the CPU engine on five of them and
  level on the patch tutorial (Section~\ref{sec:cg31}).
\item Two modelling-side levers (Section~\ref{sec:modelling}): a run-end criterion on the port probes
  (Table~\ref{tab:probeend}; the energy criterion stalls on the static charge left behind by openEMS's Gaussian pulse,
  which has a DC component, so the runs always went to their maximum time; upstream has since added an opt-in DC-free
  pulse, commit d686f1d, which is not in this work's base), and a mesh-line merging rule with a minimum cell size whose
  sweep on the SMA launch gives another 2.0$\times$ for 0.12\,dB of $S_{11}$ at $-23$\,dB (Section~\ref{sec:meshfloor}).
  Since 2026-11-04 the toolkit's defaults use the shorter settings of both (a hold of 0.5 instead of 1; on the SMA
  launch a 0.05\,mm floor instead of 0.02\,mm); together with
  the CPU engine of cg18 they cut a complete simulation of the SMA launch from 562.7\,s with stock to 37.4\,s and one
  of the GCPW line from 76.5\,s to 5.2\,s (Table~\ref{tab:e2e}), and with cg29 the Metal engine ran these runs
  1.22$\times$ and 1.37$\times$ as fast as the CPU engine (Table~\ref{tab:e2ecg29}, paired per round, on a loaded
  machine).
\item A list of things that did not work, with numbers (Section~\ref{sec:rejected}), and the artefact set
  (Section~\ref{sec:repro}).
\end{itemize}

\paragraph{Scope.} All numbers are from one machine, an Apple M5 MacBook with 16\,GB, macOS~27.0.1, Apple clang~21,
except two validation-suite runs made on a Linux virtual machine (Section~\ref{sec:repro}).
They were measured while the machine was also used for other work (the harness reports the load of other processes
with every number: sampled in the second before each run and, since the cg27 tables of 2026-11-04, their CPU and GPU
use during it,
Section~\ref{sec:method}). Stock openEMS and the fork were both built with Apple
clang, the default on macOS; upstream's FAQ~\cite{biergaizi29} advises building openEMS with GCC on macOS because
clang generates inefficient code for the SSE engine. That was measured for the inner loop only, in the standalone
kernel, which holds a copy of the stock loop and of its rewrite (one thread, three grids outside the valley;
\code{results/kernel\_gcc\_vs\_clang\_2026-11-03.txt}): GCC~16 runs the stock loop 1.16--1.39$\times$ faster than
Apple clang, and the rewrite of snapshot AB (Section~\ref{sec:engine}) at about the same speed with both, so against
a GCC-built stock the rewrite's one-thread gain in the kernel is 1.26--1.47$\times$ instead of 1.74--1.80$\times$. In
the same kernel the rewrite gives the stock loop's results bit for bit under clang but not under GCC, which contracts
it differently (Section~\ref{sec:contract}). A GCC build of openEMS itself was not made, so the factors of this paper
are against stock built with Apple clang, the default compiler on macOS, and the bit-identity checks used clang builds.
The synchronisation and hook changes do not depend on the compiler. Five details are platform-specific and guarded:
thread QoS (macOS), the starting point of the automatic thread search (the performance-core count on macOS, half the
available hardware threads elsewhere; since cg29 the flag waits on macOS spin only while the thread count is at most
that count, and other CPUs keep the barrier spin time), the separable far-field path (built only with Apple's
Accelerate; other builds keep the original loop), the opt-in flush-to-zero of subnormals (any arm64 CPU) and the
Metal diamond engine (Apple only, opt-in, built with \code{ENABLE\_METAL}). The fused sweep's mechanism, fewer
passes over the fields, is not specific to this CPU, but it was measured only here. The allocation-placement valley
depends on where the allocator puts the two field arrays and was seen only with the allocator of macOS~27.0.1 on the
M5 (Section~\ref{sec:valley}); the stagger that avoids it is applied on every platform and costs nothing measurable
outside the valley.

\section{Background}\label{sec:background}

\subsection{The EC-FDTD update and its memory traffic}
openEMS stores voltages $v$ (the E-field integrated along a cell edge) and currents $i$ (the H-field integrated along
a dual edge), three components each, as 32-bit floats, and advances them with the usual leap-frog:
\begin{align}
v_n &\leftarrow vv_n\, v_n + vi_n\,(\text{curl of } i \text{ around the edge}), \\
i_n &\leftarrow ii_n\, i_n + iv_n\,(\text{curl of } v \text{ around the dual edge}),
\end{align}
with four coefficients per component and cell. Stock openEMS already keeps the coefficients in a \emph{compressed}
operator (\code{Operator\_SSE\_Compressed} in v0.37.0): \code{f4vector}s of four cells (Section~\ref{sec:engines})
whose coefficients agree lane by lane share one entry of a table, and each \code{f4vector} carries a 32-bit index
into it, 1\,B per cell. The index array itself is allocated with one 32-bit entry per cell, of which one per vector
is used; on meshes with few z lines the unused entries share cache lines with the used ones, so the engine streams up
to the full 4\,B per cell (in stock and in the fork up to cg28; cg29 allocates one entry per z-vector,
Section~\ref{sec:cg29}). The table is small (2\,500--69\,500 entries for the models here, under 14\,MB: 1.5\,MB for
the tutorials, 13\,MB for the SMA launch; from the project log) and largely cache-resident; the per-cell traffic that
remains is the fields themselves. With separate E and H passes, each pass reads and writes its own field (12\,B each
way) and reads the other field (12\,B) from DRAM; with the index that is 37\,B per update and roughly 74 bytes per
cell and time step. This is a traffic model, not a counter measurement, as are all bytes-per-cell figures in this
paper; it counts the compact 1\,B of index, and Section~\ref{sec:engine} also gives the figures with the index as
allocated. At 150\,GB/s, a sweep that moves these 74 bytes per cell from DRAM in every time step runs at
2\,000\,\MC\ at best~\cite{roofline} (an engine that keeps a tile's fields on chip over several time steps moves fewer,
Section~\ref{sec:gpu}); the stock engine's 1\,041\,\MC\ on a 7.9-million-cell model is 77\,GB/s, half the bus.
Smaller models fit in part in a last-level cache
of a few tens of MB and could in principle run faster; stock ran them \emph{slower} (277--318\,\MC), which is the
first hint that something other than bandwidth dominated there.

\subsection{Engines, operators, extensions}\label{sec:engines}
openEMS separates the operator (coefficients and mesh, built once) from the engine (the time loop). The SSE engines
work on \code{f4vector}s, four floats in one 16-byte vector. With $n_v=\lceil N_z/4\rceil$ z-vectors per z-line, one
vector holds the cells $z$, $z+n_v$, $z+2n_v$ and $z+3n_v$; the arrays are stored component-fastest (C order
\code{[x][y][z\_v][component]}, the I-J-K-N layout of Section~\ref{sec:engine}), and the z extent is padded to $4n_v$.
Boundary conditions, excitations, lumped elements, dispersive materials and the cylindrical multigrid are
\emph{extensions}, each with up to six hooks per time step (pre-/apply-/post- for voltages and for currents). The
multithreaded engine splits the x range into one slab per thread and runs each hook of each extension between
barriers.

\subsection{The machine}
Apple M5 (2025): 4 performance (P) and 6 efficiency (E) cores, 16\,GB LPDDR5X unified memory, 153\,GB/s
nominal~\cite{m5}. The GPU has no memory of its own; it competes with the cores for the same bus. We measured the
ceiling with two standalone kernels (a $200^3$ box with perfect-electric-conductor (PEC) walls, no PML; from the
project log, not in \code{results/}): a material-table CPU loop reaches 1\,491\,\MC\ at 72 bytes per cell (107\,GB/s)
and the same kernel in Metal 1\,763\,\MC\ (127\,GB/s). The GPU gets 15--20\,\% more of the bus than the cores, but no
bandwidth of its own. Everything that follows is therefore about moving fewer bytes, fewer times, with fewer stalls.
